% ==============================================================================
% T0 Theory / FFGFT: Document 281 (Chapter Content, DE)
% Titel: Der Goldene Schnitt in FFGFT -- Notation, Rollen und Status
% Autor: Johann Pascher
% Datum: 15. Juni 2026
% Zweck: Konsolidierende Klaerung. Die frueheren Dokumente werden NICHT
%        veraendert; dieses Dokument legt die kanonische Lesart fest und
%        verweist auf alle substantiellen Vorkommen von phi im Korpus.
% Referenzen: Dok. 005, 009, 018, 019, 022, 023, 028, 035, 036, 083, 133,
%        147, 149, 150, 153, 154, 159, 172, 187, 188, 190 (P35/P37), 271,
%        272, 274, 275; Massenherleitung Dok. 006.
% ==============================================================================

\section*{Abstract}
Der Goldene Schnitt $\varphi = (1+\sqrt5)/2$ tritt im FFGFT-Korpus in
mehreren, klar zu unterscheidenden Rollen auf -- und unter zwei verschiedenen
Schriftzeichen, die teilweise mit anderen Groessen (Skalarfeld, Winkel,
Potential, Wellenfunktion) kollidieren. Dieses Dokument aendert die frueheren
Texte \emph{nicht}; es legt die kanonische Notations- und Lesart fest, ordnet
jedes substantielle Vorkommen einer Rolle zu und trennt sauber, was aus
$\varphi$ \emph{hergeleitet}, was nur \emph{behauptet} und was schlicht
\emph{falsch gerechnet} ist. Kernaussagen: (i)~$\varphi$ ist in FFGFT nirgends
Fundament -- das ist $\xi=4/30000$; (ii)~die einzige $\varphi$-Verwendung mit
echtem geometrischem Gehalt ist $\arctan(\varphi^{-3})$ in der CP-Phase
(Dok.~172), weil $\varphi^{-3}=\sqrt5-2$ die natuerliche Zwischenskala der
Fibonacci-Konvergenz ist; (iii)~$\varphi$-Massenleitern (Dok.~172/188) sind
heuristisch, numerisch fehlerhaft und fallen unter die
$\varphi^N$-Trivialitaetswarnung P35; (iv)~$1/\varphi$ ist Durchgangspunkt,
kein Fixpunkt (Dok.~275); (v)~$\pi$ und $e$ werden nirgends hergeleitet.

% ==============================================================================
\section{Anlass und Geltungsbereich}
% ==============================================================================

Bei der Pflege des Korpus zeigte sich ein Notationsproblem: dasselbe
Zeichen \verb|\phi| meint in verschiedenen Dokumenten verschiedene Dinge,
und der Goldene Schnitt selbst wird mal als $\phi$ (geschlossene Form, aus
\verb|\phi|), mal als $\varphi$ (offene Form, aus \verb|\varphi|)
geschrieben. Da die alten Dokumente in der Version v1.1.2 abgeschlossen sind
und nicht mehr nachtraeglich umgeschrieben werden, dient dieses Dokument als
\emph{verbindliche Referenz}: Wer in einem aelteren Text auf $\varphi$ oder
$\phi$ trifft, schlaegt hier nach, welche Rolle gemeint ist.

\begin{important}
\textbf{Geltungsregel.} Dieses Dokument ist normativ fuer die \emph{Lesart},
nicht fuer den \emph{Wortlaut} der alten Dokumente. Wo ein alter Text und
dieses Dokument sich widersprechen, gilt fuer den physikalischen Gehalt die
hier festgehaltene, gegen Dok.~190 (P35/P37) und Dok.~275 abgeglichene
Fassung.
\end{important}

% ==============================================================================
\section{Kanonische Notation}
% ==============================================================================

Die folgende Zuordnung ist die kanonische Konvention. Sie beschreibt zugleich,
\emph{wie} der Korpus die Zeichen faktisch verwendet -- mit der bekannten
Inkonsistenz beim Goldenen Schnitt.

\begin{center}
\begin{adjustbox}{max width=\textwidth}
\begin{tabular}{l l l l}
\hline
\textbf{Befehl} & \textbf{Glyphe} & \textbf{Kanonische Bedeutung} & \textbf{Bemerkung} \\
\hline
\verb|\varphi| & $\varphi$ (offen) & Goldener Schnitt $(1+\sqrt5)/2$ & empfohlenes Standardzeichen \\
\verb|\phi|    & $\phi$ (geschlossen) & Skalarfeld / Winkel / Index & NICHT Goldener Schnitt \\
\verb|\Phi|    & $\Phi$ (gross) & Potential / komplexes Feld & z.\,B.\ $\Phi=\rho e^{i\theta}$ \\
\verb|\psi|    & $\psi$ & Wellenfunktion & Quantenmechanik \\
\hline
\end{tabular}
\end{adjustbox}
\end{center}

\begin{caution}
\textbf{Bekannte Altlast.} In v1.1.2 schreiben mehrere Dokumente den
\emph{Goldenen Schnitt} faelschlich als \verb|\phi| statt \verb|\varphi|
(u.\,a.\ Dok.~005, 009, 022, 023, 028, 035, 036, 083, 133, 147, 187, 188).
Gleichzeitig nutzen Feldtheorie-Dokumente \verb|\phi| korrekt fuer das
\emph{Skalarfeld} (z.\,B.\ Dok.~202: $\phi_{n,l,j}$ mit
$\Box\,\phi_{n,l,j}=-m^2\phi_{n,l,j}$) und \verb|\phi| als
\emph{Wicklungsindex} $n_\phi$ in $(n_\theta,n_\phi)$ (z.\,B.\ Dok.~051, 202).
Innerhalb \emph{eines} Dokuments ist $\phi$ jeweils eindeutig; die Ueberladung
besteht nur \emph{zwischen} Dokumenten. Daher darf nicht pauschal
\verb|\phi|$\to$\verb|\varphi| ersetzt werden -- das wuerde die
Feld-/Winkel-Dokumente zerstoeren.
\end{caution}

\noindent\textbf{Numerische Festlegung.} Durchgehend gilt
\begin{equation}
  \varphi = \frac{1+\sqrt5}{2} = 1{,}6180339887\ldots, \qquad
  \varphi^2 = \varphi + 1, \qquad
  \varphi^{-1} = \varphi - 1 = 0{,}6180\ldots,
\end{equation}
\begin{equation}
  \varphi^{-3} = \sqrt5 - 2 = 0{,}2360679\ldots, \qquad
  \arctan(\varphi^{-3}) = 13{,}2825^\circ .
\end{equation}

% ==============================================================================
\section{Die Rollen von $\varphi$ im Korpus}
% ==============================================================================

Aufbauend auf der Dreiteilung in Dok.~190 (P37) -- (1)~mikroskopisch/$\xi$-Ebene,
(2)~makroskopisch/$1/\varphi$-Durchgangspunkt, (3)~Teilchenphysik -- wird hier
feiner aufgeloest. Jede Rolle ist mit ihren Fundstellen und ihrem Status
versehen.

\subsection{Rolle A -- Selbstaehnlichkeits-Attraktor (behauptet)}
\textbf{Fundstelle:} Dok.~172, \S\,,,Der goldene Schnitt als fraktaler
Attraktor``. \textbf{Aussage:} Die fraktale Feldgeometrie mit $D_f=3-\xi$
erzwinge eine Selbstaehnlichkeitsrekursion, deren einziger stabiler Fixpunkt
$\varphi$ sei (algebraisch $\varphi^2=\varphi+1$).
\textbf{Status: behauptet, nicht hergeleitet.} $\varphi^2=\varphi+1$
\emph{ist} die Definitionsgleichung von $\varphi$; der eigentliche Schritt
,,$D_f=3-\xi$ \emph{erzwingt} genau diese Rekursion`` wird assertiert, nicht
gezeigt. Die Variante in Dok.~188 (,,$\varphi$ als stabilstes inkommensurables
Verhaeltnis``) ist staerker motiviert -- $\varphi$ ist via Kettenbruch
$[1;1,1,\dots]$ die ,,irrationalste`` Zahl --, haengt aber ebenso daran, dass
die Geometrie genau dieses Verhaeltnis selektiert.

\subsection{Rolle B -- CP-Phase $\arctan(\varphi^{-3})$ (echter Gehalt)}
\textbf{Fundstelle:} Dok.~172, \S\,,,Formaler Beweis: $\arctan(\varphi^{-3})$
als geometrische Notwendigkeit``. \textbf{Aussage:} Die Fibonacci-Quotienten
$F(n{+}1)/F(n)$ konvergieren gegen $\varphi$ mit Fehlern der Ordnung
$\varphi^{-2n}/\sqrt5$; das geometrische Mittel von erster und zweiter Ordnung
ist $\sqrt{\varphi^{-2}\varphi^{-4}}=\varphi^{-3}=\sqrt5-2$. Mit der
Torsionsphase folgt
$\delta = 270^\circ + \arctan(\varphi^{-3}) \approx 283{,}28^\circ$
(NOvA/T2K 2024: $283{,}1^\circ\pm14^\circ$).
\begin{correct}
\textbf{Status: am besten begruendete $\varphi$-Verwendung des Korpus.} Hier
ist $\varphi^{-3}$ keine Wahl, sondern die natuerliche Zwischenskala der
Fibonacci-Konvergenz; der Exponent $3$ ist ueber die Z3-Ordnung
($3$~Generationen/Sektoren) vorwaerts motiviert. Das ist qualitativ
verschieden von einer angepassten $\varphi$-Potenz.
\end{correct}

\subsection{Rolle C -- Feinstrukturkonstante $\alpha=\xi\,\varphi^3\,F(7)$}
\textbf{Fundstelle:} Dok.~172, \S\,,,$\alpha$ aus reiner Geometrie``.
\textbf{Aussage:}
\begin{equation}
  \alpha \approx \xi\cdot\varphi^3\cdot F(7) = \xi\cdot\varphi^3\cdot 13
  \;\Rightarrow\; \alpha^{-1} = 136{,}19 \quad(-0{,}62\%).
\end{equation}
\textbf{Status: numerisch korrekt nachgerechnet,} aber als eine von mehreren
$\alpha$-Routen zu verstehen. Die strenge $\alpha$-Herleitung laeuft ueber
$\alpha=\xi E_0^2$ mit $E_0=\sqrt{m_e m_\mu}$ und der Potenzbeziehung
$\alpha\propto\xi^{11/2}$ (Dok.~011); die $\varphi^3\cdot13$-Form ist die
geometrisch-,,saubere``, aber gegenueber dem Messwert um $0{,}62\%$ abweichende
Variante.

\subsection{Rolle D -- Massen-/Generationenleiter (heuristisch, fehlerhaft)}
\textbf{Fundstellen:} Dok.~172 ($m_\mu/m_e\approx\varphi^{10}$),
Dok.~188 (Leptonleiter $\varphi^5\pi^2$, $\varphi^5$), Dok.~009
($\varphi^{12}/(1-\xi)^2$), Dok.~022/035/133 ($\varphi^{\text{gen}}$-Skalierung).
\textbf{Status: heuristisch; in Dok.~172/188 numerisch falsch} (siehe
Abschnitt~\ref{sec:fehler}). Die belastbare Massenherleitung erfolgt
\emph{nicht} ueber $\varphi$, sondern ueber $\xi$ und Quantenzahlen (Dok.~006,
Abschnitt~\ref{sec:masse}). Diese Rolle faellt unter P35.

\subsection{Rolle E -- Fraktaler Skalierungsfaktor $\alpha=\varphi^{-1}$
(KEIN Fehler)}
\textbf{Fundstelle:} Dok.~188, \S\,,,Fraktale Kopplungskonstante``.
\textbf{Aussage:} $\alpha=\varphi^{-1}=0{,}6180$ als Skalierungsfaktor zwischen
fraktalen Ebenen, $(x_{k+1},\dots)=\varphi^{-1}(x_k,\dots)$.
\begin{important}
\textbf{Status: eigene, klar definierte Groesse -- nicht die
Feinstrukturkonstante.} Im selben Dokument steht die Feinstruktur sauber
getrennt als $\alpha_{\text{EM}}=\xi E_0^2=1/137{,}04$. Es liegt also
\emph{kein} Rechenfehler vor, sondern eine bewusste Doppelbelegung des
Buchstabens $\alpha$. Reine Lesbarkeitsfrage; Dok.~190 P37(a) fuehrt diese
Verwendung bereits.
\end{important}

\subsection{Rolle F -- Durchgangspunkt $1/\varphi$ der $\xi$-Rekursion}
\textbf{Fundstelle:} Dok.~275. \textbf{Aussage:} Die Daempfungsfolge
$r(k{+}1)=r(k)(1-\xi)$, $r(0)=D_f/3$, passiert $1/\varphi$ bei
$k^* = \log\varphi/\xi \approx 3609$ (exakt $k^*_{\text{exakt}}=3608{,}51$).
\begin{caution}
\textbf{Status (Revision 11.\ Juni 2026): Durchgangspunkt, kein Fixpunkt.}
Einziger Fixpunkt der Rekursion ist $0$. Da $k^*(c)$ fuer \emph{jedes} Ziel
$c$ existiert, ist die Beobachtung unverankert, solange $1/\varphi$ nicht
unabhaengig motiviert wird -- die $\varphi^N$-Trivialitaetswarnung P35 greift.
\end{caution}

\subsection{Rolle G -- Pentagonale Symmetriebrechung}
\textbf{Fundstellen:} Dok.~018 (anomales $g{-}2$), Dok.~149 (Torsion),
Dok.~019, 153, 154. \textbf{Aussage:} $\varphi$ als Trager der pentagonalen
(fuenfzaehligen) Symmetriebrechung. \textbf{Status: strukturelles Motiv;}
teilt den Begruendungsstand von Rolle~A (die Selektion gerade der pentagonalen
Symmetrie ist die offene Stelle).

\subsection{Rolle H -- HLV-Substrat (nicht FFGFT)}
\textbf{Fundstelle:} Dok.~271. \textbf{Aussage:} In Hidden-Layer-Vacuum (HLV,
Krueger) ist $\varphi$ die Substrat-Geometrie. \textbf{Status: gehoert HLV,
nicht FFGFT.} Beim Bruecken-Vergleich darf $\varphi$ nicht stillschweigend von
der HLV- auf die FFGFT-Seite uebertragen werden.

\subsection{Rolle I -- Triviale Darstellung $X=\varphi^N$}
\textbf{Fundstellen:} Dok.~190 (P35), 271, 272, 274. \textbf{Aussage und
Status:} Jede Magnitude laesst sich als $\varphi^N$ \emph{oder} $\xi^N$
schreiben; ein Exponent allein ist keine Vorhersage. Gilt fuer beide Basen
gleichermassen.

% ==============================================================================
\section{Ehrliches Register: korrekt / offen / falsch}\label{sec:fehler}
% ==============================================================================

\subsection{Korrekt nachgerechnet}
\begin{itemize}
  \item $\alpha = \xi\,\varphi^3\cdot13 \Rightarrow \alpha^{-1}=136{,}19$
        ($-0{,}62\%$), Dok.~172.
  \item $\varphi^{-3}=\sqrt5-2=0{,}23607$, $\arctan(\varphi^{-3})=13{,}28^\circ$,
        Dok.~172.
  \item $1/\varphi=0{,}6180$ als Durchgangspunkt bei $k^*\approx3609$,
        $r(3609)=0{,}617994\ldots$, Dok.~275 (Status: unverankert).
  \item $245 = 5\cdot7^2 = \varphi^{12}/(1-\xi)^2 \approx 244{,}98$, Dok.~009 --
        mit \emph{deklariertem} empirischem Anker $1{,}314$ (P10/P37(d)).
\end{itemize}

\subsection{Numerisch falsch -- in den Altdokumenten, hier korrigiert}
\begin{critical}
\textbf{Dok.~188, Leptonleiter (Z.~$\approx$318):}
\begin{align}
  \text{behauptet: } & \frac{m_\mu}{m_e}\approx\varphi^5\cdot\pi^2\approx 206{,}7
    && \text{tatsaechlich } \varphi^5\pi^2 = 109{,}46 \;(\text{gemessen }206{,}768),\\
  \text{behauptet: } & \frac{m_\tau}{m_\mu}\approx\varphi^5\approx 16{,}94
    && \text{tatsaechlich } \varphi^5 = 11{,}09 \;(\text{gemessen }16{,}82).
\end{align}
\textbf{Dok.~172 (Z.~$\approx$356):} $m_\mu/m_e\approx\varphi^{10}=122{,}99$ --
$40\%$ daneben und im Widerspruch zur (ebenfalls falschen) Dok.-188-Form.
Selbst arithmetisch korrigiert braeuchte man krumme Exponenten
($\varphi^{11{,}08}$ bzw.\ $\varphi^{5{,}87}$): es gibt keine glatte
$\varphi$-Potenz, die die Leptonverhaeltnisse trifft.
\textbf{Konsequenz:} Die $\varphi$-Leptonleiter ist als Numerologie zu
markieren (P35-Nullbefund), nicht als geometrische Konsequenz.
\end{critical}

\subsection{Die belastbare Leptonmassen-Herleitung -- ueber $\xi$, nicht
$\varphi$}\label{sec:masse}
Dok.~006 setzt $m_i = r_i\,\xi^{p_i}\,v$ mit aus Quantenzahlen bestimmten
Exponenten:
\begin{equation}
  p_e = \tfrac32,\; r_e\approx 1{,}349; \qquad
  p_\mu = 1,\; r_\mu = \tfrac{16}{5} = 3{,}2;
\end{equation}
\begin{equation}
  \frac{m_\mu}{m_e} = \frac{r_\mu}{r_e}\,\xi^{p_\mu-p_e}
  = 2{,}372\cdot\xi^{-1/2} = 2{,}372\cdot86{,}6 = 205{,}4 \quad(-0{,}65\%).
\end{equation}
Der \emph{Exponentenunterschied} $-\tfrac12$ ist der strukturelle, vorwaerts
gerichtete Gehalt; die $O(1)$-Koeffizienten ($r_e$ ist rueckgerechnet) tragen
den Restfit (vgl.\ P40: nur Verhaeltnisse exakt). $\varphi$ wird hier nicht
benoetigt und fuegt nichts hinzu.

% ==============================================================================
\section{Was NICHT aus $\varphi$ (oder $\xi$) folgt}
% ==============================================================================

\begin{itemize}
  \item \textbf{$\pi$ wird nicht hergeleitet.} $\pi$ ist die vorausgesetzte
        Kreis-/Torus-Konstante. Die schoenste Deutung (Dok.~159): auf dem Kreis
        konvergiert $\sum 1/n^2=\pi^2/6$, die periodische Geometrie erzwingt
        Konvergenz -- aber $\pi$ ist gedeutet, nicht abgeleitet.
  \item \textbf{$e$ (Eulersche Zahl) wird nicht hergeleitet}, sie tritt nur als
        $\exp(\cdot)$ auf.
  \item \textbf{Emergent abgeleitet} sind dagegen: $13=F(7)$ (Fibonacci) und
        $27=3^3$ (Z3-Struktur, drei Generationen).
\end{itemize}

% ==============================================================================
\section{Vorkommens-Tabelle (substantielle Faelle)}
% ==============================================================================

Die Tabelle listet die \emph{inhaltlich relevanten} Vorkommen. Reine
Definitions- oder Durchlaufnennungen folgen derselben kanonischen Lesart.

\begin{center}
\begin{adjustbox}{max width=\textwidth}
\begin{tabular}{l l l l}
\hline
\textbf{Dok.} & \textbf{Zeichen} & \textbf{Rolle (s.\ o.)} & \textbf{Status} \\
\hline
172 & $\varphi$ & A Attraktor, B CP-Phase, C $\alpha$, D $\varphi^{10}$ & B korrekt; A behauptet; D falsch \\
188 & $\phi$ & D Leptonleiter, E Skalierungsfaktor & D falsch; E ok (eigene Groesse) \\
275 & $\varphi$ & F Durchgangspunkt $1/\varphi$ & korrekt, aber unverankert (P35) \\
190 & $\varphi$ & I Trivialitaet P35, Rollen-Map P37 & normativ \\
009 & $\phi$ & D $\varphi^{12}/1{,}314$ & empirischer Anker deklariert \\
271 & $\varphi$ & H HLV-Substrat, I Trivialitaet & Abgrenzung HLV/FFGFT \\
272 & $\varphi$ & I $X=\varphi^N$ keine Vorhersage & korrekt \\
274 & $\varphi$ & I Trivialitaet gilt fuer $\varphi$ und $\xi$ & korrekt \\
018 & $\varphi$ & G Pentagonale Symmetriebrechung ($g{-}2$) & strukturelles Motiv \\
149 & $\varphi$ & G Pentagonale Symmetriebrechung (Torsion) & strukturelles Motiv \\
019, 153, 154 & $\varphi$ & G Pentagonale Symmetriebrechung & strukturelles Motiv \\
022, 035, 133 & $\phi$ & D $\varphi^{\text{gen}}$-Skalierung & heuristisch (P35) \\
036 & $\phi$ & D kosmische Skalen $r_n=r_0\varphi^n$ (BAO) & heuristisch (P35) \\
083, 147, 187 & $\phi$ & D Qubit-/Photonik-Frequenzen & heuristisch \\
005, 028, 035, 150 & $\phi$/$\varphi$ & Definition $(1+\sqrt5)/2$ & Grundlage \\
241, 261 & $\varphi$ & F ,,zur Ruhe gekommene`` Rekursion & vgl.\ Dok.~275 \\
\hline
202 & $\phi$ & \emph{Skalarfeld} $\phi_{n,l,j}$, Index $n_\phi$ & NICHT Goldener Schnitt \\
051 & $\phi$ & \emph{Wicklungsindex} $n_\phi$ & NICHT Goldener Schnitt \\
067, 097, 129 & $\phi$ & \emph{Skalarfeld} (Lagrange/QFT) & NICHT Goldener Schnitt \\
\hline
\end{tabular}
\end{adjustbox}
\end{center}

% ==============================================================================
\section{Empfehlung}
% ==============================================================================

\begin{enumerate}
  \item \textbf{Notation kuenftig:} Goldener Schnitt ausschliesslich
        \verb|\varphi| ($\varphi$). \verb|\phi| bleibt Feld/Winkel,
        \verb|\Phi| Potential, \verb|\psi| Wellenfunktion.
  \item \textbf{Altdokumente:} nicht umschreiben; diese Referenz gilt.
  \item \textbf{$\varphi$-Massenleitern (Dok.~172/188):} als P35-Numerologie
        lesen, nicht als Herleitung. Belastbar ist allein die $\xi$-Route
        (Dok.~006).
  \item \textbf{Echter $\varphi$-Gehalt} bleibt: $\arctan(\varphi^{-3})$ in der
        CP-Phase (Dok.~172) und die Begriffstrennung der drei $\varphi$-Rollen
        (Dok.~190 P37).
\end{enumerate}
